\section{Formalism of the Grassmann tensor}
\label{sec:formulation}

The variables $\eta_i \ (i=1,\cdots,N) $ are single component Grassmann numbers 
which satisfy the anti-commutation relation 
$ \{\eta_{i},\eta_{j} \}=0 $.
We begin with defining a Grassmann tensor and its contraction rule with single component index $\eta_i$. 
Then those with multi component index $\Psi=(\eta_1,\eta_2,\cdots,\eta_N)$ are defined by extending the single 
component case straightforwardly. 


The {\it Grassmann tensor} $\mathcal{T}$ of rank $N$ is defined as
\begin{align}
\label{eq:def1}
	\mathcal{T}_{\eta_{1}\eta_{2}\cdots\eta_{N}}=\sum_{i_{1}=0}^1 \sum_{i_{2}=0}^1 \cdots \sum_{i_{N}=0}^1 T_{i_{1}i_{2}\cdots i_{N}}\eta_{1}^{i_{1}}\eta_{2}^{i_{2}}\cdots\eta_{N}^{i_{N}},
\end{align}
where %a set of Grassmann numbers $\eta_i$, $i=1,\cdots,N$, satisfying $\{\eta_{i},\eta_{j}\}=0$, 
$\eta_i$ are single component Grassmann numbers
and $T_{i_{1}i_{2}\cdots i_{N}}$ is referred to as a coefficient tensor, 
whose rank is also $N$, with complex entries. 
Fig.~\ref{fig:fig1}~(a) represents a Grassmann tensor, 
where the external lines correspond to the indices $\eta_i$.  


We consider a {\it Grassmann contraction} among Grassmann tensors. 
Let $\mathcal{A}_{\eta_1\ldots\eta_N}$ and $\mathcal{B}_{\zeta_1\ldots\zeta_M}$ 
be two Grassmann tensors of rank $N$ and $M$, respectively.
%
\footnote{ We assume that either of $\mathcal{A}$ and $\mathcal{B}$ 
is a commutative tensor whose coefficient tensor $T_{i_1i_2 \cdots i_N}=0$ 
for $(i_1+i_2 +\cdots i_N) \, {\rm mod} \, 2 =1$. 
In this case, 
$\mathcal{A}_{\eta_1\ldots\eta_N} \mathcal{B}_{\zeta_1\ldots\zeta_M}
=\mathcal{B}_{\zeta_1\ldots\zeta_M} \mathcal{A}_{\eta_1\ldots\eta_N}$ 
and Eq.~(\ref{eq:oriented}) can also be expressed as
$
	\int\diff\bar{\xi}\diff\xi~\e^{-\bar{\xi}\xi}
\mathcal{B}_{\bar{\xi}\zeta_2 \ldots \zeta_M}  \mathcal{A}_{\xi\eta_2\ldots \eta_N}. 
$
%
%
%since the order of $\mathcal{A}$  and  $\mathcal{B}$ 
%does not affect the result.
}
%
The Grassmann contraction has an orientation
 which comes from the anti-commutation relation of 
Grassmann variables.  
We define a Grassmann contraction from $\eta_1$ to $\zeta_1$ as
\begin{align}
\label{eq:oriented}
	\int\diff\bar{\xi}\diff\xi~\e^{-\bar{\xi}\xi}
%\mathcal{A}_{\psi\eta}\mathcal{B}_{\bar{\eta}\phi}.
\mathcal{A}_{\xi\eta_2\ldots \eta_N}\mathcal{B}_{\bar{\xi}\zeta_2 \ldots \zeta_M}.
\end{align}
Eq.~(\ref{eq:oriented}) itself is a Grassmann tensor, and
the coefficient tensor of Eq.(\ref{eq:oriented}) is 
given by a contraction of two coefficient tensors  of $\mathcal{A}$ and $\mathcal{B}$ 
with some sign factors. 
One can consider a contraction from $\eta_i$ to $\zeta_j$ 
as a straightforward extension of Eq.(\ref{eq:oriented}) 
with keeping the weight factor $\e^{-\bar{\xi}\xi}$.
In Fig.~\ref{fig:fig1}~(b), the Grassmann contraction is shown 
as a shared link with the arrow.
Note that $\int\diff\bar{\xi}\diff\xi \, \e^{-\bar{\xi}\xi}
\mathcal{A}_{\bar\xi\eta_2\ldots \eta_N}\mathcal{B}_{{\xi}\zeta_2 \ldots \zeta_M}$ 
should be represented as  Fig.~\ref{fig:fig1}~(b) with the opposite arrow. 



\begin{figure}[htbp]
\centering\includegraphics*[width=0.7\hsize]{graph_single.pdf}
\caption{Graphical representations of 
(a) Grassmann tensor Eq.(\ref{eq:def1}) and  (b) Grassmann contraction Eq.~(\ref{eq:oriented}). 
The external lines specify uncontracted indices.  
The arrow in the internal line represents a contracted direction from $\xi$ to $\bar \xi$. 
}
\label{fig:fig1}
\end{figure}

Let us now move on to the multi component case. For simplicity of explanation, 
we take $N=mK$ for Eq.~(\ref{eq:def1}). Then $N$ Grassmann numbers $\eta_n$ 
are divided into $m$ component variables $\Psi_a$ ($a=1,\cdots,K$) 
as $\Psi_a=(\eta_{(a-1)m+1}, \eta_{(a-1)m+2} ,\cdots,\eta_{am})$. 
The Grassmann tensor Eq.~(\ref{eq:def1}) is also expressed as 
\begin{align}
\label{eq:def_multi}
	\mathcal{T}_{\Psi_1 \Psi_2 \cdots \Psi_K}  \equiv  \mathcal{T}_{\eta_{1}\eta_{2}\cdots\eta_{N}}
\end{align}
The rank of a Grassmann tensor should be carefully read from the dimension of indices 
since both sides of Eq.~\eqref{eq:def_multi} have the same rank.
Fig.~\ref{fig:fig1m}~(a) represents an example of Grassmann tensor with multi component indices, 
where the external links shown as solid lines correspond to the multi component indices $\Psi_a$. 
Other cases are straightforwardly generalized from 
Eq.(\ref{eq:def_multi}) which is the case of $N=4m$ and $m$ component indices.


\begin{figure}[htbp]
\centering\includegraphics*[width=0.7\hsize]{graph_multi.pdf}
\caption{Graphical representations of
(a) Grassmann tensor Eq.(\ref{eq:def_multi}) and  (b) Grassmann contraction 
Eq.~(\ref{eq:oriented_multi}). 
The external lines specify uncontracted multi component indices.  
The arrow in the internal line represents a contracted direction from $\Xi$ to $\bar \Xi$. 
}
\label{fig:fig1m}
\end{figure}

To define the Grassmann contraction with multi dimensional indices, 
we consider the case of $N=mK, M=mL$ in Eq.(\ref{eq:oriented}) 
for simplicity.  
The $N$ and $M$ rank tensors $\mathcal{A}_{\eta_1\eta_2\cdots\eta_{N}}$ and 
$\mathcal{B}_{\zeta_1 \zeta_2\cdots \zeta_{M}}$ are expressed as 
$\mathcal{A}_{\Psi_1 \Psi_2 \cdots \Psi_K}$ and $\mathcal{B}_{\Phi_1 \Phi_2 \cdots \Phi_L}$ 
where $\Psi_a$
and $\Phi_a$ are $m$ component indices defined as in Eq.~(\ref{eq:def_multi}). 
Then the Grassmann contraction is given for the multi component case: 
\begin{align}
\label{eq:oriented_multi}
	\int\diff\bar{\Xi}\diff\Xi~\e^{-\bar{\Xi}\Xi}
%\mathcal{A}_{\psi\eta}\mathcal{B}_{\bar{\eta}\phi}.
\mathcal{A}_{\Xi\Psi_2\ldots \Psi_K}\mathcal{B}_{\bar{\Xi}\Phi_2 \ldots \Phi_L}
\end{align}
where $\Xi=(\xi_1,\xi_2,\cdots,\xi_m), 
\bar \Xi=(\bar \xi_m,\cdots,\bar \xi_2, \bar \xi_1)$ and 
\begin{align}
	\diff\bar{\Xi}\diff\Xi~\e^{-\bar{\Xi}\Xi} \equiv 
\prod_{n=1}^{m}\diff\bar{\xi}_{n}\diff\xi_{n}~\e^{-\bar{\xi}_{n}\xi_{n}}.
\end{align}
The case of $m=1$ reproduces Eq.~(\ref{eq:oriented}). We should note that 
$\bar \Xi$ contains $\bar\xi_n$ in a reverse order so that
the coefficient tensor of Eq.(\ref{eq:oriented_multi}) is simply given by 
a tensor contraction of coefficient tensors of $\mathcal{A}$ and $\mathcal{B}$
without extra sign factors. 
Fig.~\ref{fig:fig1m}~(b) shows the Grassmann contraction with multi component indices.  


It is easy to define a {\it Grassmann tensor network} with these notations. 
Let $\mathcal{T}_n$ be Grassmann tensors. Then 
the tensor network is defined by a product of 
$\mathcal{T}_1\mathcal{T}_2\cdots$ where all indices 
are contracted as  Eq.~(\ref{eq:oriented_multi}).



